% Part: lambda-calculus
% Chapter: introduction
% Section: unique-readability

\documentclass[../../../include/open-logic-section]{subfiles}

\begin{document}

\olfileid[tr]{lam}{syn}{unq}
\olsection{Tek Okunabilirlik}

Her terimin yalnızca tek bir biçimde oluşturulup oluşturulamayacağını
sorabiliriz; gerçekten de öyledir. Her lambda teriminin tek bir kuruluşu ve
tek bir çözümlemesi vardır. Aşağıdaki tartışmada \emph{kuruluş},
\olref[trm]{defn:term} içindeki kuruluş kurallarını bir ya da birkaç kez
kullanarak bir terim oluşturma işlemidir.

\begin{lem}\ollabel{lem:term-start}
Bir terim ya bir değişkenle ya da bir parantezle başlar.
\end{lem}

\begin{proof}
Bir ifade ancak \olref[trm]{defn:term} uyarınca kurulmuşsa terim sayılır.
\olref[trm]{defn:term-var} sonucundaysa bir değişken olmak zorundadır.
\olref[trm]{defn:term-abs} ya da \olref[trm]{defn:term-app} sonucundaysa bir
parantezle başlar.
\end{proof}

\begin{lem}\ollabel{lem:app-start}
Bir uygulamanın sonucu ya iki parantezle ya da bir parantez ve bir değişkenle
başlar.
\end{lem}

\begin{proof}
$M$ bir uygulamanın sonucuysa $(PQ)$ biçimindedir ve dolayısıyla bir
parantezle başlar. $P$ bir terim olduğundan, \olref{lem:term-start} gereği ya
bir parantezle ya da bir değişkenle başlar.
\end{proof}

\begin{lem}\ollabel{lem:initial}
Bir terimin hiçbir öz başlangıç parçası kendi başına bir terim değildir.
\end{lem}

\begin{prob}
\olref[lam][syn][unq]{lem:initial} lemasını terimlerin uzunluğu üzerinde
tümevarımla ispatlayın.
\end{prob}

\begin{prop}[Tek Okunabilirlik] \ollabel{prop:unq}
Her terimin tam olarak bir kuruluşu vardır.
\end{prop}

\begin{proof}
  Bunu terimlerin kuruluşu üzerinde tümevarımla ispatlarız.
  \begin{enumerate}
    \item $M$, bir $x$ değişkeni için $x$ biçiminde olsun. Soyutlama ve
      uygulamaların sonuçları her zaman parantezle başladığından bunlar
      $M$'yi kurmak için kullanılmış olamaz. Dolayısıyla $M$'nin kuruluşu,
      \olref[trm]{defn:term}\olref[trm]{defn:term-var} kuralının tek bir
      uygulaması olmak zorundadır.
    \item $M$, bir $x$ değişkeni ve bir $N$ terimi için
      $(\lambd[x][N])$ biçiminde olsun. Tek bir değişken olmadığından
      \olref[trm]{defn:term}\olref[trm]{defn:term-var} uyarınca kurulmuş
      olamaz. \olref{lem:app-start} gereği bir uygulamanın sonucu da değildir.
      O hâlde $M$, yalnızca $N$ üzerinde bir soyutlamanın sonucu olabilir.
      Tümevarım varsayımına göre $N$'nin kuruluşu da tektir.
    \item $M$, $P$ ve $Q$ terimleri için $(PQ)$ biçiminde olsun. Bir
      parantezle başladığından
      \olref[trm]{defn:term}\olref[trm]{defn:term-var} ile kurulmuş olamaz.
      \olref{lem:term-start} gereği $P$, $\lambd$ ile başlayamaz; dolayısıyla
      $(PQ)$ bir soyutlamanın sonucu olamaz. Şimdi $M$'nin uygulamayla başka
      bir kuruluşu bulunduğunu, örneğin aynı zamanda $(P'Q')$ biçiminde
      olduğunu varsayalım. Bu durumda $P$, $P'$nin öz başlangıç parçası ya da
      $P'$, $P$'nin öz başlangıç parçası olurdu; bu ise
      \olref{lem:initial} ile olanaksızdır. Dolayısıyla $P$ ve $Q$ tek türlü
      belirlenir; tümevarım varsayımına göre bunların kuruluşları da tektir.
  \end{enumerate}
\end{proof}

Yukarıdaki önermenin daha okunaklı bir anlatımı şöyledir:
\begin{prop}
  Bir $M$ terimi yalnızca aşağıdaki biçimlerden birinde olabilir:
  \begin{enumerate}
    \item $M$ tarafından tek türlü belirlenen bir $x$ değişkeni için $x$.
    \item Hem $x$ hem de $N$, $M$ tarafından tek türlü belirlenmek üzere,
      bir $x$ değişkeni ve başka bir $N$ terimi için $(\lambd[x][N])$.
    \item Her ikisi de $M$ tarafından tek türlü belirlenen iki $P$ ve $Q$
      terimi için $(PQ)$.
  \end{enumerate}
\end{prop}

\end{document}
